\section{Contributions}\label{sec:contributions}

The first contribution is the Motivo language model.
It provides a compact set of abstractions for notes, rests, chords, sequences, patterns, tracks, voices, and control
flow.
These abstractions are deliberately close to symbolic musical structure and deliberately far from raw MIDI byte
construction.
The language is small, but it is sufficient to express repeated arrangements and parallel lines in the current
implementation.

The second contribution is the compiler architecture.
The implementation uses a conventional staged pipeline: scanner, parser, semantic analyzer, lowerer, and backend.
The most project-specific part is the cursor-based lowering model.
That model lets sequential statements advance musical time, explicit offsets place events without changing the main
cursor, voices run with isolated cursors, and pattern calls return reconstructed musical values.
This is the core mechanism that connects source-level structure to MIDI note events.

The third contribution is Motivo Studio.
The application shows how the compiler can be made accessible through a web interface.
The editor, diagnostics loop, MIDI parsing, piano-roll visualization, playback controls, server wrapper, Docker stack,
and CI workflows demonstrate the software engineering side of the project.
For a bachelor's thesis, this part is important because it shows that the work is not only theoretical.
It is an implemented system with several interacting modules.

The same contribution can also be viewed through the structural abstractions implemented by the compiler. The
language provides constructs such as \texttt{pattern}, \texttt{track}, and \texttt{voice}, which allow
polyphonic arrangements to be described without manually calculating absolute timestamps or delta-times. The
deterministic lowering algorithm connects those abstractions to a static intermediate representation, while
the semantic analyzer checks type and structural constraints before MIDI generation.

The fourth contribution is the empirical evaluation.
The local benchmark results show that Motivo can generate large symbolic event streams from compact source programs and
compile them quickly in the tested Release configuration.
The results are intentionally bounded by the current implementation and hardware.
They support the practical claim that Motivo is suitable for interactive local experimentation at the current scale.
